\subsection*{Note (3)}
The third note, on pp.~53--54, extends the program attached to \S3,
p.~16, beyond the varieties called anabelian. Grothendieck proposes a
description of the category of schemes of finite type over the absolute
base through systems of profinite groups expressing their étale topoi,
and then envisages passage to more general schemes by limits. His phrase
``qui fait semblant'' marks the presentation's deliberate avoidance of
rings and equations while retaining the schemes' étale topoi. The
explicit qualification, introduced as ``un grain de sel'', is that the
reconstruction is to concern the category obtained by making universal
homeomorphisms invertible; in the finite-type setting under discussion
he describes these morphisms as finite, radicial, and surjective. This
places the proposed dictionary of \S3 in the language of categories and
localization and supplies further context for the priorities listed in
note (4). The closing paragraph on p.~54 presents this ambition as an
outcome of Galois theory and answers the familiar complaints about its
purpose and generality. The note enlarges and qualifies the proposed
reconstruction; it does not supply a completed construction of that
dictionary. The related archive title \emph{La Longue Marche à travers
la théorie de Galois} is already named with \S3.
(\href{https://grothendieck.umontpellier.fr/web/viewer.html?file=../119.pdf\#page=56}{Copy 1, pp.~53--54, beginning at site 55};
\href{https://grothendieck.umontpellier.fr/web/viewer.html?file=../119.pdf\#page=57}{closing paragraph at site 56};
\href{https://grothendieck.umontpellier.fr/web/viewer.html?file=../119.pdf\#page=19}{call in \S3, p.~16, site 18}.)
